\chapter{黑格尔《法哲学原理》中的传统与革命}

1833年5月，黑格尔逝世一年半以后，爱德华·甘斯出版了收入逝者友人版《黑格尔全集》中的《法哲学原理》第2版，并以对此书“未来命运”预兆般的暗示结束了他热情洋溢的前言：他发现“整部著作系由一块自由的金石所铸就”。作为黑格尔体系的一部分，甘斯说道，它将与整个体系共沉浮。也许，若干年后，它将与体系一道进入人们的观念、进入普遍的时代意识，届时它的“术语”将会消失，它的“深邃”将成为共同财富。但是，甘斯进一步总结道，其后，这种在此著中以及一般而言通过概念把握其时代的哲学将在哲学上过时，而仅具有历史意义；新的哲学发展已经在“已然改变的现实”中做好了准备：“这样，它的时代在哲学上已告终结，它成为历史的一部分。一种出自同一基本原理的新的哲学发展脱颖而出，以达到对现实变化的一种不同的理解。”\footnote{《全集》第8卷，前言，柏林，1833年，第VII页。}甘斯此言在不少方面都是预言性的。他道出此言时是黑格尔学派唯一的且依然非常虔诚的造反派。这段话预言了1830年代晚期和40年代黑格尔学派为了争夺这个体系的合法遗产而进行的争斗，也预言了——当然会令人想起1830年七月革命以及黑格尔对它的保留态度——1840年后现实的改变，最后还道出了黑格尔法哲学的特殊命运：在黑格尔法哲学出现后不久，它的时代就会在哲学上成为过去，然而它的历史影响才刚刚开始。甘斯趁“此书”出版之机所言的“未来命运”很快就在整个黑格尔体系上面得到应验。老年黑格尔主义、青年黑格尔主义和新黑格尔主义对法哲学的独特分析引人注目；事实上，这本书在1833年即已成为历史，而在随后的年代（首先是在洛伦茨·冯·施泰因和马克思这里）还“创造了”历史。

与《精神现象学》《逻辑学》《美学》《哲学史》和《历史哲学》在19世纪和20世纪所产生的巨大历史影响不同，《法哲学原理》和《自然哲学》的命运相同：它们受到同时代人的激烈排斥，随着实定的自然科学和历史法学的发展而陷入矛盾之中。对黑格尔法哲学的流俗成见——它迎合了普鲁士国家和1820年代的复辟时期——在一定程度上导致了从斯塔尔(Stahl)到埃德曼、海姆、卢格、罗森克兰茨、费舍尔(Fischer)这些批判者和评论家的空疏，1930年左右新黑格尔主义的哲学努力[首先可以在拉伦茨(Larenz)、杜尔凯特(Dulckeit)、布塞(Busse)和晚期宾德尔(Binder)那里见到]也是如此。\footnote{新黑格尔主义对《法哲学》的评论，见科勒(J. Kohler)发表在《法哲学和经济哲学档案》(Archiv für Recht- und Wirtschaftsphilosophie)第V卷(1912年)上面的论文；合集：宾德尔、布塞、拉伦茨《黑格尔法哲学导引》(Einführung in Hegels Rechtsphilosophie, 1931)；宾德尔(J. Binder)：《法哲学奠基》(Grundlegung zur Rechtsphilosophie, 1935)以及K. 拉伦茨和G. 杜尔凯特的作品，法哲学的新黑格尔主义开端为人们认为在黑格尔那里作为“国家中的人民共同体”而重新发现的“具体的共同体自身”[参见G. 杜尔凯特：《法概念和法形态》(Rechtsbegriff und Rechtsgestalt, 1936)，第17页]。黑格尔法哲学在19世纪和20世纪的影响史尚有待书写。}这里再度证实了甘斯所写前言的惊人准确性：前言当时就指出“此书的实际价值与对它的接受与传播极不相称。”\footnote{《全集》第8卷，前言，柏林，1833年，第V页。}众所周知，1840年后的政治运动进一步损害了对黑格尔法哲学的承认。它对时代的影响，不是来自对其意义的恰当评论，而是来自对它的批判。因此，法哲学的时代在哲学上已经终结，因为1840年后的时代不再是哲学的，而是政治的、经济的和历史的了。

毫无疑问，在法和国家领域，19世纪黑格尔的伟大辩护人非洛伦茨·冯·施泰因莫属。但是在他的《当代法国的社会主义和共产主义》（1842年）及其国家和社会理论的相同概念背后，它们所描述的社会和历史现实已经发生了多么巨大的变化啊！拉萨尔自称是黑格尔和甘斯的真正传人；就其哲学才能和法学教养而言，在黑格尔的学生中，他是1860年后唯一能够在法哲学基础上写出一部“政治学”的人。然而非常典型的是，他在他的第一部同时也是最后一部法哲学著述《既得权利的体系》（1861年）中宁愿投身于一个特殊的法学问题，向当时的经验主义和历史主义致敬。马克思觉得自己既不是黑格尔的辩护士，也不是他的传人（在这个词语的真正意义上），而是他的批判性的战胜者；因此，他就是在他的时代唯一如其所是地对待1821年《法哲学原理》的人：将它作为“唯一与正式的当代现实保持在同等水平上的德国历史”。\footnote{《〈黑格尔法哲学批判〉导言》，载《马克思恩格斯全集》（第1版），第一部分，第1卷第1册（1927年），第612页。（中译文引自《马克思恩格斯全集》第3卷，人民出版社2002年版，第205页。——译者）}青年马克思于1843年对第261-313节写下的《黑格尔国家法批判》（马克思通过这部著作彻底摆脱了经由革命改变了的政治形而上学传统），是19世纪黑格尔传承中对法哲学所做的唯一的、有时在历史方面还相当深刻的、与黑格尔的论述水平相当的评论。

显然，这部令人遗憾的未完成的评论性批判作品并未提出黑格尔法哲学在政治哲学传统中的地位问题，按其历史意图，它根本就不可能提出这个问题。与甘斯的预言一致，马克思的评论认为黑格尔法哲学仅仅属于历史，因此他的批判一开始就已经预设，法哲学的时代在哲学上已然过去。众所周知，批判的结果就是使哲学脱离法，而转向经济学。同时，这种公然宣称的非哲学依然深受这部最后的伟大的法的形而上学的影响，因为马克思不是在世界这部大书中，而是首先在黑格尔法哲学中学会看见并理解到了作为现代之特征的“对立”：国家与社会、国民与私人、政治生活与社会生活、中世纪与现代之间的对立。接下来的解释努力也恰好在其中见到了黑格尔的“深刻”：如马克思的评论所言，黑格尔“处处从各种规定（如我们各个邦所存在的那类规定）的对立开始，并且强调这种对立”。\footnote{《马克思恩格斯全集》（第1版）第一部分，第1卷第1册（1927年），第465页。[中译文引自《马克思恩格斯全集》第3卷，第69-70页。——译者]}对《法哲学原理》（首先是对其第三部分）的这种哲学与历史上的问题和形式的分析表明，这种对立一直深入形式的结构之中，并且整部著作阐述了一部在革命的基础上生长起来的非亚里士多德“政治学”的方案，其原理试图描绘和规定现代人的基础——他们的国家、社会、家庭和个人，同时又不割断与古典传统的联系。

\section*{1}

黑格尔给他的这部著作起了两个书名：《法哲学原理或自然法和国家学纲要》，这种做法乍看起来，不免令人侧目。第二个书名并非如人们所料，是一个副名；相反，通过添加“纲要”(im Grundrisse)二字，而成为法的“原理”(Grundlinien)的异名。由此，对“法哲学”形式上的强调便具有一种深意。

这一异名中的两个名称指出了黑格尔以前形而上学思想中的两门学科，即自然法和国家学。自然法是在近代欧洲（尤其是在17、18世纪）发展起来的，国家学则以“政治学”之名，属于欧洲的古老传统，并在直到沃尔夫为止的学院哲学中地位稳固。然而，旧政治学的本质恰好在于，其中并无自然法和国家学的分离。因为对于直到17世纪为止的欧洲思想（在德国则要推迟到18世纪中叶）来说，政治学是对人、公民或者政治的社会（或者共同体）在法律上进行规治的共同体进行全面阐述的理论，亚里士多德对于城邦作为政治共同体的论述乃是政治学的典范。\footnote{对此，参见布伦纳：《社会史的新道路》第30、207页，还有我在不久之后对市民社会所做的概念史研究。}由于城邦或国家作为政治共同体或者市民社会包含了自由市民（即基于其身份而享有政治权利的市民）在自由和法下的全部生活，因此非政治的和市民以外的人就表现为野蛮的，而“市民的”(civilis)一词便通过其与“自然的”(naturalis)的对立而获得了积极的含义。只有在16、17世纪，从马基雅维利、博丹和霍布斯开始，当现代国家从政治上组织起来的市民社会中解放出来时，国家学才从旧政治学中解放出来，从而18世纪出现了自然法和国家学之间的对立。此后，市民的(civilis)和自然的(naturalis)的意义便完全颠倒了过来。因为从历史—政治角度，而非按照其哲学内容来看，自然法表达了这样一种努力，即要么坚决限制现代国家对其本身原本就拥有的良好法秩序的社会的（旧的意义上的）“不合法的”干预，要么用一种新的自由和法的概念革命性地克服已经僵化为单纯的权力功能的国家政治。自然法与国家学之间的对立（或者如德国人所言，道德和政治之间的对立）在精神上引导了近代欧洲革命，并在此后与革命过程共进退。黑格尔法哲学以这种对立为前提，并在革命的历史背景下，如革命本身一样，想要通过哲学政治的努力达到对革命的克服。黑格尔在此著前面写上两个书名，作为其书名，正好表明了这一点。

“法哲学原理”是这部著作的第一书名，也是它的真正名称。我们看到，这个书名在第二个书名（“自然法和国家学”）前面。作为对现代国家的阐述，黑格尔的国家学想要成为一种突出意义上的“法哲学”。因此，在他这里，政治学，作为国家中人的社会生活的理论，便成了法哲学，因为他试图通过将前国家的自然法与（政治的权力功能为个人颁布的）法完全建立在自然—理性法——人之为人的自由——之上，扬弃二者之间的革命性对立。然而，对黑格尔来说，自由的这种自然—理性法（这里不详述其意义）具有一种历史内容。自从基督教进入世界历史起，这种内容就是灵魂在上帝面前的平等，从18世纪末开始，就是个人在革命国家面前的平等和自由。由于黑格尔理解到了革命的历史意义，并将革命的历史内容——所有人的自由和权利能力——接纳到他的政治思想中，因此对他来说，权利便在其与自由的联系中成为政治哲学的中心。随着革命的这一结果，革命以前的自然法和国家学之间的对立便失去了意义。此后，政治学的基础就只能是权利，但不是旧的市民社会的具有权利能力的市民的权利，而是权利能力的权利、追求其自由人的“概念”。

然而，黑格尔一方面想要在一种建立在权利之上的政治哲学中积极调和自然法和国家学之间的对立，另一方面他又回到了革命前的政治学传统。在这一传统中，无此种“对立”。如沃尔夫翻译这个名称时所言，政治学是关于“人的社会生活（特别是共同体）”的理论，此外别无他物。\footnote{克里斯蒂安·沃尔夫的德语版“政治学”全名为：《克里斯蒂安·沃尔夫致热爱真理者，为促进人类幸福对人的社会生活特别是共同体的理性思考》(Vernünftige Gedanken von dem gesellschaftlichen Leben der Menschen und insonderheit der gemeinen Wesen zur Beförderung der Glückseligkeit des menschlichen Geschlechts, den Liebhabern der Wahrheit mitgetheilet von Christian Wolff)，1721年第1版。}因此，说黑格尔法哲学的前提建立在普遍的法概念和自由要求之上（这两个成分表明它隶属于近代自然法），这是正确的：它作为国家理论（并且在国家理论中）发现它的基础。如同在亚里士多德那里，城邦是“真正的基础”，按照现实性来说是“首要的东西”，按照理性来说是人联合成社会的一切结合的“目的”。\footnote{黑格尔：《法哲学原理》，E. 甘斯出版(1833年)，第312页及以下。下面的引用均出自这一版本。}事实上，法哲学在很大程度上就是国家学，众所周知，从埃德曼和海姆、拉萨尔和罗森茨威格直到拉伦茨和杜尔凯特的新黑格尔主义对它的解释都集中于这一点。国家学这一名称必然令人想到旧政治学传统：在国家整体中阐述一族人民或者一个社会的政治秩序，黑格尔有意识地重拾这一传统。与古老的欧洲政治学传统相契合，黑格尔拒绝时人喜爱的“关于国家的哲学”及其目标：提出自己的“理论”，“世界上从未有过国家和国家制度，现在也还没有，但是现在——这个现在是永远继续下去的——似乎应该从头开始”。\footnote{同上书，第7页。（中译文引自《法哲学原理》，第4页。——译者）}他像传统一样，与当前的伦理、国家及其秩序紧密相连。同时，革命性的时代政治理论已经打破了传统意识，学院哲学以及康德的政治著作中非同寻常的辩护性论争清楚表明了这一点。这方面尤其典型的是革命后对“哲学与现实间的关系”的特别反思。国家的现实性及其实际的法和统治的形式既不是（如《法哲学原理》序言中所言）通过“建立彼岸物”而以乌托邦的方式摧毁掉，也不是在传统的政治或市民的科学或历史(scientia bzw. historia politica sive civilis)的意义上达到对其关系的简单陈述。相反，黑格尔的政治思想与现实的“肯定”关系意味着通过“理解当前和现实的东西”而“探究理性的东西”，因此本身就是一种现实与概念的“关系”，并且是一种“关于”国家与历史的哲学（如果人们想要违背黑格尔的话）。看似恢复传统的东西，实际上是一种通过革命破坏了的、其意义已然改变的传统。黑格尔的国家学正是在现代自然法的基础上生长起来，并且将要表明，政治在法中的中介，是一个既进行中介也在运动、逃离中介的中心。——由此，我们远远超出了与黑格尔法哲学的两个书名有关的法、自然法和国家学的概念分析，进至法哲学本身的体系问题及其与政治哲学传统的关系。

\section*{2}

可以假定，1820年黑格尔不仅在考虑到给他的柏林听众发一本讲座（“我循职进行法哲学授课”）提纲的需要时，认为有必要对《哲学全书》的相关部分进行单独阐述。\footnote{《法哲学原理》，第3页。（中译文参见《法哲学原理》，第1页。——译者）}相反，如果比较1817年海德堡版《哲学全书》“客观精神”相关章节与1821年《法哲学原理》，就会清楚地看到，它并非如柏林版序言所言，只是“对同一些基本概念所作的更为详尽尤其是更加系统的阐述”，而且还讨论了更多的东西。而通过附释澄清“正文的抽象内容”以及“更广泛地考虑当前流行的一些见解”的提示，也不能作为应当付印《法哲学原理》的“动因”。决定性的可能更是这一事实，即1817年后哲学体系的发展导致黑格尔超出《逻辑学》，进到“法的科学”。因为他在序言中说：“我已在我的《逻辑学》中详尽阐述了思辨知识的本性。”\footnote{同上书，第4页。（中译文参见《法哲学原理》，第2页。——译者）}就其阐述了“实践”知识的本质而言，《法哲学原理》正好与前者匹配。如黑格尔所强调的那样，法哲学正好关乎“科学”，尽管它并未总是在逻辑形式上非常清楚地显明“对象的具体、自身如此多样的属性”。实际上，在黑格尔本人生前出版的著作中，《法哲学原理》（1821年）与《逻辑学》（1812、1816年）彼此并列。如果将《精神现象学》视为体系（在1817年《哲学科学百科全书》中，这个体系总称为知识的总体）的导论，那么逻辑学和法哲学的体系撰构上的显著平行便会令人想到哲学的传统划分：理论哲学和实践哲学。这种划分存在于欧洲形而上学中，在18世纪，沃尔夫及其学派也在对亚里士多德学术传统的依赖中系统地坚持了这一划分。这里，实践哲学与逻辑学并列，人类灵魂的两方面的能力(facultas aminae)——认识能力和欲求能力(facultas cognoscitiva atque appetitiva)——被称为两者的基础(fundamentum)。自然法、政治学、伦理学的划分基于实践哲学的定义：“真正哲学的教导运用欲求能力求善避恶的这一部分，叫实践哲学。因此实践哲学就是指导欲求能力求善避恶的科学。”(Ea vero philosophiae pars, quae usum facultatis appetitivae in eligendo bono et fugiendo malo inculcat, philosophia practica dicitur. Est adeo philosophia practica scientia dirigendi facultatem appetitivam in eligendo bono et fugiendo malo.)\footnote{克里斯蒂安·沃尔夫：《理性哲学或逻辑学》（1740年），哲学概论之预备性讨论，第3章，第62节。}

特别的是，黑格尔没有在他的著作中援引过这个传统，因此需要考虑他在柏林时期对大学“完善的哲学课程”的意见。黑格尔认为，“完善的哲学课程”与明确、中肯、系统的阐述一样严肃和重要。这里的“大致”意思是说，“提供一种按照理论哲学讲座与实践哲学讲座的习惯划分而来的完善性。因此，我在一个学期讲授逻辑和形而上学，在同年度的另一学期讲授自然法和国家学，或是包含了伦理学或义务论意义上的法哲学”。\footnote{《柏林大学哲学系档案》，文献目录第2号第1卷，载《纽伦堡著作集》，J. 荷夫迈斯特出版，1938年，第XXIV页。}如《逻辑学》第1版序言中最尖锐地表达出来的那样，他确信旧形而上学已经终结。这一确信说明了黑格尔在哲学的古老划分上面已然断裂了的传统意识，因为前述引文不过是“大致”的回忆。公开表达出来的对于“大约25年来哲学思维方式在我们中间所遭受的完全改变”的意识、“精神的自我意识在这个时期对它自身所达到的立场”并不是偶然地导致他再次改变了他于1807年计划的《科学体系》的体系结构，并且放弃了精神现象学为一方面，逻辑学、自然哲学和精神哲学为另一方面的两分法。此后，黑格尔的哲学体系划分为他所言的“两门实在科学”，即自然哲学和精神哲学，逻辑学作为“纯粹科学”构成了它们的第一部分。\footnote{见《逻辑学》，第一部分，拉松编，1934年，序言，第7-8页。[中译文参见〔德〕黑格尔：《逻辑学》（上卷），杨一之译，商务印书馆1966年版，第5-6页。——译者]}

这里不能进一步讨论这个划分的问题。对我们的主题来说，根本在于，众所周知，法哲学与历史哲学一道从属于“精神哲学”第二部分即“客观精神”之下。客观精神的内容是人类世界的活动(Pragmata)：法和国家、经济和社会、政治和历史。因此，客观精神学说即以人与人之间，在家庭、市民社会和国家中有其基础和目的的那些人类行为为主题。这些行为就其本身的意义而言并不超越于人世之外，而是人世的“客观精神”，在欧洲形而上学传统中，它们是实践哲学（实践哲学首先是政治哲学和伦理学，因而是广义上的道德形而上学）的对象。与其对象的制度、在社会和政治上有序的人类世界相一致，直到18世纪为止，基于古代文化传统的学院哲学在实践哲学领域按照实践对象本身的结构和人的特定身份，区分了三门学科：伦理学、家政学和政治学。因为如沃尔夫所言，人的身份可以用两种方式来考察——“一个人要么在一定程度上是人，要么在一定程度上是市民。或者换而言之，一个人要么在一定程度上生活在人类社会或自然状态，要么在一定程度上生活在市民社会。因此实践哲学便以两种方式分为两个部分。”(vel quatenus est homo, vel quatenus est civis, aut, quod perinde est, vel quatenus vivit in societate generis humani seu statu naturali, vel quatenus vivit in societate civili. Ob duplicem hunc respectum philosophia practica in duas dispicitur partes.)\footnote{克里斯蒂安·沃尔夫：《理性哲学或逻辑学》(1740年)，哲学概论之预备性讨论，第3章，第63节。}这两个部分就是伦理学和政治学。伦理学的对象为人类的伦理状态，因此，它研究一般的道德存在者(entia moralia)和人之为人的地位。“对生活在自然状态或人类社会中的人进行思考的哲学部门，叫伦理学。因此，我们将伦理学定义为指导自然状态下（或者人就是他自己的主人、不服从别人的权力时）自由行动的科学。”(ea philosophiae pars, in qua homo consideratur tamquam vivens in statu naturali, seu in societate generis humani, Ethica appellatur. Quamobrem Ethicam difinimus per scientiam dirigendi actiones liberas in statu naturali, seu quatenus sui juris est homo, nulli alterius potestati subjectus.) 政治学的对象为市民社会(societas civilis)的法制度，因此不是类的伦理自然，而是通过法律规定的国家(res publica、Staat)的政治自然。它的主题不是作为人的人，而是作为市民的人，因此将他从他的伦理行为的普遍性和无规定性中拉出来，置于国家的市民法律的特殊性之下。“对生活在国家或者市民状态下的人进行思考的哲学部门，叫政治学。因此，政治学就是指导市民社会或者国家中自由行动的科学。”(Ea philosophiae pars, in qua homo consideratur tamquam vivens in Rep. seu statu civili, Politica vocatur. Est itaque Politica scientia dirigendi actiones liberas in societate civili seu Republica.)\footnote{克里斯蒂安·沃尔夫：《理性哲学或逻辑学》(1740年)，哲学概论之预备性讨论，第3章，第64、65节。}

在实践哲学这两门基本学科（政治学作为市民及其在市民社会中的权利和义务的理论，伦理学作为整个人类社会中的人的理论）以外，还有将人作为家庭成员来考察的家政学。如同亚里士多德《政治学》卷一，沃尔夫谈到了构成家庭之基础的“小型社会”(societates minores)：“先于国家或市民社会，有处于自然状态或者国家以外的小型社会，如夫妻社会、父权社会、主奴社会，它们通常来自同一家族。”(Dantur praeter Remp. seu societatem civilem societates aliae minores, quibus etiam in statu naturali seu extra Remp. fuisset locus, veluti societas conjugialis, paterna, herilis, quae simplicium nomine venire solent.)\footnote{克里斯蒂安·沃尔夫：《理性哲学或逻辑学》(1740年)，哲学概论之预备性讨论，第3章，第66节。}

在这个持续了好几百年的传统中，实践哲学的这种三分法不只是僵死的学院知识。它表达了革命前欧洲社会状况，反映了其政治、经济和伦理的风貌。随着17、18世纪革命，现代社会从国家中解放出来，这个传统瓦解了：先是在英格兰（霍布斯在政治学上是反亚里士多德派，洛克、弗格森和休谟已经完全站在新“社会”的基础上），然后在法国。伦理学与政治学分离，作为“道德哲学”，变成了它的敌人，家政学突破了其特有的限制，变成了18世纪的国家经济学和国民经济学——如黑格尔所言，变成了“在现代世界的基础上产生的若干科学之一”。\footnote{《法哲学原理》，第189节，第255页。（中译文参见《法哲学原理》，第204页。——译者）}在现代社会从国家中解放出来的过程中，以及与此相关的实践哲学（特别是旧政治学）的内容变化中，家政学中的革命发挥了决定性作用。对此，我们将在黑格尔这里具体指出来。有别于英格兰和法兰西，唯有在德国，18世纪仍然突出地保持了实践哲学的形而上学传统。这里，在沃尔夫及其学派，这一传统还在描绘着革命前的社会：家政学和政治学对应于家庭社会和市民社会的不同领域、国家和社会的统一，伦理与有效的法律相连，因此伦理学与政治学相连。\footnote{对此，参见青年洪堡对1785-1786年通俗哲学家恩格尔按照沃尔夫学派手册讲授的实践哲学的部分的笔记[《洪堡全集》(Gesammelte Schriften)，第一部分，第7卷第2分册，1908年，第363页以下]。}

对于当时西欧思想来说具有建构性的自然法与政治学的分离，在18世纪德国学院哲学中（除了大法学家普芬道夫和托马修斯以外）并未发挥决定性作用，这一点绝非偶然，因为自然法脱离了传统实践哲学的框架。西欧近代自然法的功能是与政府对立，保护“社会”的权利，使社会从政府解放出来，与此相对，沃尔夫将自然法的内容限定为一种普遍的实践哲学“理论”。自然法的定义在于它的规定之中：善行和恶行的科学，出现在欲望之中的人对善恶的内部知识。沃尔夫总结道：“自然法显然是实践哲学（即伦理学、政治学和家政学）的理论。因此，区分理论和实践并不难，自然法本身就能传达伦理学、家政学和政治学。”(Facile patet jus naturae esse theoriam philosophiae practicae, Ethicae scilicet, Politicae atque Oeconomicae. Quamobrem cum non opus sit, theoram a praxi distingui; jus naturae in ipsa Ethica, Oeconomica atque Politica tradi potest.)\footnote{克里斯蒂安·沃尔夫：《理性哲学或逻辑学》(1740年)，哲学概论之预备性讨论，第3章，第68节，补充。}沃尔夫只是形式上将自然法与实践哲学分开，它的内容则能够“传达”到社会的伦理、经济和政治的生活秩序之中去。

\section*{3}

现在，如果比较一下沃尔夫依赖亚里士多德和经院哲学最后一次在其庞大著作中建立起来的这个欧洲传统实践哲学的体系纲要与黑格尔政治形而上学的外部轮廓，仔细考察后就会惊奇地发现，尽管政治革命和18世纪下半叶开始的工业革命粉碎了传统实践哲学体系的整个结构，但是它的基础在黑格尔的《法哲学原理》中依旧宛然可见。概念的自我展开作为法哲学的辩证奠基（在黑格尔看来，在实践中，概念就是在其自由中的人的意志，自由就是意志的“概念”），导致他将这个领域划分为“抽象法”“道德”和“伦理”三个部分。抽象法和道德包含了人的前政治存在的理论。在抽象法中，可以认识到演绎为——以财产自由和职业自由为顶点的一——现代市民的私人权利的自然法；在道德中，可以认识到那种沉浸于主体的自我感之中的伦理学，通过18世纪晚期革命，形成了这种主体的自我感。在抽象法和道德中，人作为人，从政治状态和政治制度下解放出来，实现自己，表达自己：或是在与外部的关系中，与外部世界和周围世界的关系中，在所有权（第41-71节）、契约（第72-80节）和“不法”（第82-104节）等形式中，另一方面在与内部的关系中，在故意和责任（第105-118节）、意图和福利（第119-128节）、善和良心（第129-141节）等形式中。这里不可能对根本而言为探讨私法内容的18世纪自然法与黑格尔的抽象法（特别是私法思想的形而上学基础方面）进行具体比较。\footnote{首先参见《法哲学原理》的详细导言，黑格尔在意志的形而上学理论的基础上阐述了现代法权人格及其普遍自由要求的“概念”（第1-32节）。在黑格尔看来，不是“以一定身份而得到考察的人”作为不同于奴隶、农奴或“侨民”的“市民”，而是“人格性本身”构成法和所有权概念的基础（第40节）。}我们仅限于指出，对黑格尔来说，自然法与伦理学两者，都是人类世界伦理—政治状态的抽象物：作为私法的自然法，标明为“抽象法”，是抽象人格性的法；18世纪主观化了的伦理学，作为道德，是同样抽象的主观性的道德。作为人的人只有在作为市民的人与作为家庭、社会和国家之内的伦理—政治的存在中，才“在实践上”实现其自身。

相比较而言，抽象法和道德探讨了传统政治哲学所不知道的现代内容——想想人在外部世界的实现（第41节以下）、在劳动和使用中对人格的所有权要求的奠基（第56、59节以下）、所有权和法权人格的等同（第51节）以及人格的主体存在、作为“道德的概念”的“意志对它本身的内部关系”——，伦理则明确回到旧的欧洲伦常传统的立场。在黑格尔这里，伦理指的是个人与特定人民和国家的“伦理力量”（第145页）及“必然关系”（第148节）的同一、古代市民的“习俗”(Ethos)。伦理学和政治学同样地包含于这种习俗之中，以保存法、伦理和国家的古老统一。黑格尔在其讨论实践哲学主题的第一篇文章《论自然法的科学探讨方式》（1802年）中写道：“个人的伦理，以及反过来，个人伦理的本质绝对是实在的（因而就是普遍的）绝对伦理。个人伦理是整个体系的脉动，并且自身就是整个体系。这里我们注意到语言的暗示。这种暗示尽管在以前遭到人们的拒绝，但从前面所言来看，它是完全正当的。这就是：绝对伦理的自然中即有一种普遍物或者伦常，因此，表示‘伦理’的希腊词，还有德语词，都卓越地表明了伦理的这种自然，而最近的伦理体系由于以一种自为存在和个别性为原则，就无法做到使用这些词语而不言明它们的联系。”语言的这种内在暗示如此重要，以至于18世纪的体系为了标明“其本质”，即个别化了的个人，“就不能滥用这些语词，而是采用了‘道德’一词。尽管从其词源上看，这个词语也指向了同样的方向，但因为它是一个新造词，就不能如此直接地反对其蹩脚的含义”。\footnote{见《批判的哲学杂志》第2卷第3册(1803年)，第1-2页。}

文献中反复提到黑格尔国家概念的“古代化”手法，同时毫无疑问，古代城邦理想对于他拒绝将个人从伦理力量的束缚中革命性地解放出来，对于他将国家制度和自由意识、伦理和法等同起来，是多么重要。因此，就形式而言，法哲学第三篇与古代实践哲学第三部分相同：“伦理”代表了伦理学和政治学的统一，对于伦理法的国家制度的传统理论——国家或者市民社会(civitas sive societas civilis)的理论——来说，这种统一是根本性的。黑格尔在他自己的“政治哲学”中并不是徒然谈到“客观的”而非“道德的”“伦理义务论”与“将在接下来的第三部分系统阐述的伦理必然性圆圈”（即国家中的“市民”关系）之间的同一性。\footnote{《法哲学原理》，第148节。一种“内在、一贯的义务论”（即旧义上的伦理学）“不外乎阐述由于自由的理念而是必然的，因而在其在国家的整个范围内是现实的那些关系”（第213-214节）。（中译文参见《法哲学原理》，第167页。——译者）}

由于黑格尔重拾古代传统，因此将现代政治和工业革命的内容包含于这种对伦理概念的古代理解之中，也就更加令人惊异。因为黑格尔《法哲学原理》的吊诡正好在于，它在此名义下记录了革命时代的基本趋势，并且在政治哲学的体系中做出了自博丹的主权概念和卢梭的公意以来最深刻的改变。唯有在这里，青年黑格尔广泛的国民经济学研究和政治研究才引人注目地获得了体系性成果。在耶拿时期的文章和讲座中，有——乍看之下令人咋舌的一一从研究英国和法国的著作获得的对于革命过程中把握到的欧洲工业和政治世界的制度的洞见，这种洞见根本上与古代文化传统无关，但与之并列。资产者被搬到了罗马帝国时期，社会的经济领域、需要、劳动和占有，被称作与希腊政治活动(politeuein)的绝对伦理相对立的“相对伦理”，古代悲剧在现代相对伦理中上演，法哲学所认识到的对“市民社会的劳动方式”的分析要么在个人与自然相关的实践过程框架下进行，要么尴尬地出现在“人民”“政府”“国家”等标题之下。\footnote{参见《批判的哲学杂志》第2卷第2册(1802年)，第79页。收入《政治学与法哲学文集》，G. 拉松编，第418页以下、464页以下，特别是参见1803-1804年《耶拿讲座》（《实在哲学》第1卷，J. 荷夫迈斯特编，1932年，第220页以下和236页以下）和1805-1806年《耶拿讲座》（《实在哲学》第2卷，1931年，第197页以下、213页以下和231页以下）。}实践哲学的旧结构再也不能容纳18世纪末政治、社会和经济的变化，这首先可以在1803-1804年和1805-1806年耶拿讲座的两份草案中看出来。这里，裂变后的世界的丰富内容并没有得到合理的体系化处理。另一方面，当黑格尔试图在1817年《哲学全书》中第一次对精神的“客观”内容进行严格体系化时，古式风格仍然使他无法对现代关系做出恰当的阐述。\footnote{唯一阐述这一主题的《海德堡哲学全书》第433节只谈到“普遍作品”及其在“诸等级”中的特殊化，其中家庭作为“个别性等级”出现。“特殊的”等级——“其作品是特殊定在的需要，而其切近的目的是特殊的主体性，但是要达到这一目的则以其他一切人的劳动为前提，同时也会影响到这一劳动”——还不是“市民社会”而只是“等级”。比较一下黑格尔在《法哲学原理》出版后在《哲学全书》第2版(1827年)对这一部分所做的全面改动（第513节以下）。在这里和其他地方，黑格尔将“伦理”术语历史化了，并将其与近代家庭、市民社会和国家的发展联系起来（第552节）。}只有1821年法哲学才能将革命的结果纳入实践政治哲学的构造，但因此也就同时摧毁了其传统的体系构架。

首先，当黑格尔将伦理，即人类世界的伦理—政治制度，它的“现实精神”，理解为家庭、市民社会和国家的著名三分时，他似乎回到了古代的欧洲传统。因为家庭和国家乃是大革命以前的政治哲学从亚里士多德（《政治学》卷第1卷，第1-2章；《尼各马可伦理学》第8卷，第14章，1162a17以下；《欧德谟伦理学》第7卷，第10章，1242a23以下）那里继受而来的主题，二者构成了人类共同生活的基础，其中国家(polis或civitas sive res publica)与社会等同，被称作“市民社会”(koinonia politike或societas civilis)，以区别于家庭的“家庭社会”(societas domestica)。\footnote{与“社会的”与“政治的”（或者“市民的”与“国家的”）之间的现代区分相对立，古代政治哲学的基本区分为“oeconomicus”（经济的）或“domesticus”（家庭的）与“civilis”（市民的）之间的区分，因此市民社会表现为公共—政治社会。}事实上，外部结构上的相似令人惊异，并且《法哲学原理》中黑格尔思想的历史来源最令人信服地表现在该书最后面的这个最重要部分对这三个概念的形式接受上。另一方面，真正说来，只有理解到这三个概念获得了完全不同的意义、这三个词语获得了全新的内容时，才能敏锐、清楚地把握黑格尔与传统的断裂。因为黑格尔法哲学所言的家庭在结构上与古老的家庭团体大相径庭，市民社会也完全不同于旧欧洲的市民社会，因此黑格尔的国家也就不同于传统的国家模式(Civitas-Modell)，因为它已经从自身中区分出“市民社会的坚实体系”。\footnote{《历史哲学》讲座导言中的这一说法，参见《历史中的理性》，J. 荷夫迈斯特编，1955年，第209页。}《法哲学原理》第三部分结构的伟大正在于，通过其思想的历史实质（即黑格尔将现代革命的各种趋势接受到法哲学之中），黑格尔对这些概念的历史来源进行了改造，并与之断裂。此后，在黑格尔伦理纲要中的传统与革命之携手同行，就应当在三个概念的结构中得到具体表达；在这三个概念中，国家和市民社会便在它们的对立中成为现代政治哲学的真正中心。

\section*{4}

黑格尔在第157节以一个简要的预先规定开始其详细的阐述。“现实精神”、人类世界伦理的实际本质是：

\begin{quote}
第一，直接的或自然的伦理精神——家庭；这种实体性向前推移，丧失其统一性，进而分裂，过渡到相对物的立场，因此就成为——

第二，市民社会，这是各个成员作为独立的个人在一个形式普遍性之中的联合，这种联合是通过成员的需要，通过作为保障人格和财产手段的法制度，以及通过维护他们的特殊利益与共同利益的外部秩序而建立的。这个外部国家——

第三，返回并凝聚为国家制度，它是实体性的普遍物的目的和现实性，也是致力于该普遍物的公共生活的目的和现实性。
\end{quote}

这一描述的引人注目之处在于，“市民社会”在这里即已占据相对较大的篇幅。家庭和国家代表并保存了旧的意义上的伦理、它们的“实体性”，市民社会则是“家庭统一性的丧失”。黑格尔引入这个概念的方式表明他显然摆脱了传统伦常，而市民社会的“矛盾”则使他在进行新的描述时对内容做出了详尽的规定。

老年黑格尔派甘斯、魏塞、H. F. W. 欣里希斯(Hinrichs)、亨宁(Leopold von Henning)已经认识到[黑格尔]将（现代意义上的）“市民社会”引入政治哲学乃是一个重大成就。毫无疑问，只有到了1820年，黑格尔才随着这个概念“理解”到了欧洲社会政治制度内部的这些基本改变，而不只是像他开始政治研究时那样进行断言。从亚里士多德到康德的法哲学传统将国家称作市民社会，因为人类社会本身就是在政治上（在自由市民的权利能力和在等级特权中）和在社会上（在家庭的经济上的重要地位）受到规治的。黑格尔则将国家的政治领域从已经变成“市民的”社会领域中区分出来。由此，“市民的”(bürgerlich)便获得了排他性的“社会的”含义，不再如18世纪那样作为“政治的”同义词使用：它指的是在绝对国家私人化了的现代市民——作为布尔乔亚的市民——的单纯“社会”地位。第190节补充用法语术语说明了这一点。这个从其政治—法上的含义解放出来的市民概念与同样解放出来的社会联系起来。它们的政治实体消解于它们两者通过现代革命而获得的社会功能之中。如此看来，黑格尔将“市民的”和“社会”结合为一个政治范畴绝非偶然。这个概念前所未有，在他自己1820年以前的著作中也未曾见到。因此，罗森茨威格所说的黑格尔对他所知道的舒尔策的《一切科学简论》（1759年）和弗格森的《文明社会史论》（1767年）的摘抄的提示就是错误的，因为此二人涉及的都是旧的传统概念（尽管形式上有所减弱）。\footnote{罗森茨威格：《黑格尔与国家》第2卷，1920年，第118页。在舒尔策这里，罗森克兰茨提到的地方在旧的意义上考察市民社会，如他所言，一个“市民国家”的社会制度（《一切科学简论》[1759年]，第189-190页）。弗格森的《文明社会史》（1767年）同样在与“政治社会”相同的意义上使用这个概念（第三部分，第6章），但在此名下除了广泛讨论政治关系（首先是斯巴达、雅典和罗马）外，还讨论了“艺术和科学”。}黑格尔从未在他的著作中接受过出现于从托马斯·阿奎那和大阿尔伯特到博丹、霍布斯、洛克和康德之中，并将革命以前的市民社会的制度表达出来的政治哲学的传统公式\footnote{参见“‘市民社会’概念及其历史起源问题”，本书第141页以下。}：“civitas sive societas civilis sive res publica latius sic dicta”（国家或者市民社会或者广义上的国家），一般而言，他也不再理会它的这种意思，就如法哲学似乎提到——从新的市民社会的立场出发——这种等同的一段话所言：“如果把国家想象为各个不同的人的统一，亦即仅仅是共同性的统一，其所想象的只是指市民社会的规定而言。许多现代的国家法学者都不能对国家提出除此以外任何其他看法。”\footnote{《法哲学原理》，第182节，补充。（中译文引自《法哲学原理》，第197页。——译者）}这里，与传统的断裂经典地表达了出来，然而另一方面，黑格尔确实完全知道旧市民社会与政治国家的同一，而且有意思的是，他是从语言出发，并且基于当时大体上还是受到等级束缚的社会[而获得这种知识的]。“尽管在这些所谓的理论看来，市民社会的一般等级和政治意义上的等级是绝不相同的东西，但语言仍然保持了以前就存在的二者之间的结合。”\footnote{同上书，第303节，第398页。（中译文引自《法哲学原理》，第323页。——译者）}同样，《法哲学原理》“市民社会”部分第一段就从实际出发、从由政治制度中解放出来的社会（即前述意义上的“市民社会”）的私人—市民出发，清楚表达了对旧概念的传统的完全抛弃：“具体个人作为特殊的人本身就是目的。他作为各种需要的整体以及自然必然性和任意的混合体，构成市民社会的一个原则，——但这个特殊的个人本质上是与其他的这样的特殊性相关的，从而每个人都通过他人的中介，并且同时完全只是作为经过了另一个原则（即普遍性的形式）的中介，使自己有效并得到满足。”\footnote{《法哲学原理》，第182节，第246页。对“市民社会”的最好阐述仍然是由罗森茨威格（《黑格尔与国家》第2卷，1920年，第118页以下）和卡尔·洛维特：《从黑格尔到尼采》（1941年）第二部分第1章第2节做出来的。此外，也可以比较克罗纳(R. Kroner)发表在《应用社会学档案》(Archiv für angewandte Soziologie)第4卷(1931年)上面的文章和里特尔的著作《黑格尔与法国大革命》对市民社会作为现代“劳动社会”的思辨阐述。}这里，获得解放的西欧社会的功利原则及其经济模式、狄德罗和爱尔维修的个人利益、边沁和富兰克林的“自利”便以德国哲学的形式出现了。于是，结果也完全一样：在个人与个人之间，产生了一个独立的、基于利益的关系网络。黑格尔在第183节写道：“利己的目的在其实现中……建立起一个在所有方面相互依赖的体系，单个人的生活和福利及其权利的定在，都同所有人的生活、福利和权利交织在一起，建立在所有人的生活、福利和权利的基础上，并且唯有在这种联系中，单个人的生活和福利及其权利的定在才是现实的和有保障的。”个体，“作为这种国家的市民，就是以其自身的利益为目的的私人”（第187节），基于这个“市民社会”特有的“必然性”（第186节），他们使自己成为其强有力的联系之“链上的一个环节”。

这个对新社会来说是典型的必然性(1)出现在其最低级阶段，即作为自然（“a. 需要及其满足的方式”）、劳动（“b. 劳动的方式”）和占有（“c. 财富”）的需要的体系（第189-208节）。这些成分以它们的方式在“社会化”的形式中产生了一种自身的必然性：无限多样的需要及其满足的手段（第190-195节）；劳动分工，“人的依赖和相互关系”变成了“完全的必然性”（第198节），同时劳动“机械化”，人可以“走开”，机器登场（第198节）；社会生产的“普遍财富”分配为特殊财富及其结果，“人的设定起来的不平等”（在黑格尔看来，“在市民社会中不但不扬弃不平等，反而从精神中产生它”[第200节]），以及社会的“全部联系”区分为各等级（第201节以下）。这种必然性也(2)支配了市民社会的上面的层次，在这一层次，在“司法”中，市民社会似乎回到了以前的制度要素，即法（第209-229节）。然而，这种司法之法本质上是私法，即适用于“市民社会中无限零星和复杂的所有权和契约的关系和种类的素材”（第213节）的“市民”法。这个领域的必然性在于：(1)“法的设定的存在”（第212节以下）；(2)其内容上的即无限多的案例以及特定的个别案例的量的特性，这些个案加剧了实定性即其作为设定的存在的特性（第214节）；(3)私人间的相互关系的无限变化的性质，这种性质限定了法，并将其限制于法律。只有第三阶段，“警察和同业公会”（第230-256节），才能将必然性“伦理化”，这种必然性像冰冷的古代命运一样支配着市民社会的个人。尽管“警察”——这是古代社会政治制度的一个残余概念——局限于外部“秩序的范围”以及平衡“无意识的必然性”（第231节以下、236节以下），但在市民工商业的同业公会，“伦理物作为内在物回到了市民社会中”（第253节以下）。

与设定起来的法相比，警察作为“普遍物的保证力量”，更加必然地同社会的特殊性不停地斗争。在即将进入国家的政治领域之前，黑格尔在“需要的体系”中转向现代社会经济的自然基础之后，便在这一部分第243-248节将工业革命的趋势纳入市民社会的概念中。此后，工业即作为“市民社会的劳动组织”而被谈及（第251节）。市民社会再也不像在沃尔夫和政治传统中那样，由“许多家庭”组成，\footnote{克里斯蒂安·沃尔夫：《克里斯蒂安·沃尔夫致热爱真理者，为促进人类幸福对人的社会生活特别是共同体的理性思考》，第214节。}而是（如果在其“不可阻挡的活动”中去理解的话）由全新的成分构成：富裕和贫困、工业和无产阶级大众、人口增长和殖民扩张。市民社会(1)“在它自身”带来了“快速发展的人口和工业”，(2)导致“大量群众下降到一定生活水平之下”，产生了所谓的“贱民”，(3)陷入一种“辩证法”，这种辩证法表现在这一事实中：“尽管财富过剩，市民社会总是不够富足。也就是说，它所占有而属于它所有的财产如果用来防止过分贫困和贱民的产生，总是不够的。”\footnote{《法哲学原理》，第245节。“国家容纳市民社会”，原文为“In-Korporation der Bürgerlichen Gesellschaft vom Staate her”，即“从国家出发吸纳市民社会”，意为在承认市民社会相对独立的前提下，将其纳入国家之下，以消除前述市民社会的“辩证法”。——译者}在黑格尔看来，由这种辩证法，对市民社会来说，只有两条出路：要么是其内部的一条内在的出路，市民社会因其矛盾趋势，而走出自身，进行殖民（第248节）。这条道路只是拖延，并不解决问题。要么是另一条——人们可能称之为——“超越的”道路，即《法哲学原理》第三部分（第257节以后）以之作结的国家容纳市民社会。然而，如1840年后的事件所表明的，这条道路并未扬弃前述辩证法，而是跳过了它。如果扩张达到了可以利用的空间的极限，殖民的道路就会对社会内部的矛盾力量无能为力，而当国家面对时代的可能性而无所作为时，就会作为乌托邦而在时代的现实性面前撞得粉身碎骨。

无论如何，黑格尔在《法哲学原理》的这个部分（第243-248节）和《历史哲学讲演录》的一个地方具体指明了市民社会是一个时代概念。黑格尔在他关于亚里士多德政治学的演讲结束时说道，“我们近代国家的抽象权利把个人孤立起来，允许他如此”，继而建立起一种必然的联系，以至于“真正说来，没有人具有整体的意识，或者为了整体而活动；他为整体而工作，但是不知道他是怎样为它工作的，他只关心他个人的保护”。\footnote{《哲学史讲演录》第2卷，K. L. 米希勒(Michelet)编，1833年，《全集》第14卷，第400页。（中译文参见《哲学史讲演录》第二卷，第364页。——译者）}黑格尔在这里用近在眼前的现代工厂形象说明现代个人、国家和社会的结构划分：“这是一种分工的活动，在这种活动中，每个人只占一份；正如在一个工厂里面，没有人自己单独制造一件完整的产品，每个人都只制造产品的一部分，不懂得制造其他各部分的技能，只有少数几个人才把各部分装配成一件产品。”只有自由的人民才有整体的意识并为它而活动，现代人作为个人本身就是不自由的——“市民的自由就是不需要普遍的事物，就是孤立的原则”（强调字体为黑格尔所加）。

《法哲学原理》第305-308节（在“伦理”第三部分，关于国家的章节）更加明确地规定了这里明言的这种时代上的联系。这几节一方面在里尔之前30年就论述了市民社会中流动的和不流动的社会力量，即等级要素的两个方面：建立在不动的自然基础之上的贵族等级和在工商业中努力拼搏的市民等级。黑格尔认为，市民等级构成了“市民社会的流动方面”（第308节）。另一方面，在将社会等级提升至国家时，这几节也将当时的两院制（世袭的贵族院和民选的代议院）固定了下来：“私人等级在立法权的等级要素中获得政治意义和政治效能。所以，这种私人等级既不是简单的不可分解的集合体，也不是分裂为许多原子的群体，而只能是它现在这个样子，就是说，它分裂为两个等级：一个等级建立在实体性的关系上，另一个等级则建立在特殊需要和以这些需要为中介的劳动上。”\footnote{《法哲学原理》，第303节，第397页。（中译文引自《法哲学原理》，第322页。——译者）}再有意思的是，黑格尔这里在用一个可以说是“政治上”修正了的市民的等级社会概念，去反对社会可能消解为“简单的不可分解的原子”，以及“许多原子的群体”，同时也反对国家的政治形式化。就等级架构上而言，市民社会依然是政治社会，并且正好处于解放了的孤立个人与国家组织之间。由此，在黑格尔这里，市民社会就不仅仅是在英国国民经济学推动下对国家学所做的一个补充，而是正好处于这两个过程的交合处：一是旧市民社会的去政治化过程，二是国家从市民社会中解放出来。因此《法哲学原理》第308节便谈到了“市民社会流动部分”的议员：“由于这些议员是市民社会派出来的，所以就直接得出这样一个结论：市民社会在派议员时是以它的本来面目出现的，就是说，它不是一群原子式地分散的单个人，不是仅仅为了完成单一的和临时的活动才一时凑合起来而事后则没有任何进一步联系的人们；相反地，它是一个分为许多有组织的协会、自治团体、同业公会的整体，这些团体因此才获得政治联系。”\footnote{《法哲学原理》，第308节，第400页以下。（中译文引自《法哲学原理》，第325-326页。——译者）}黑格尔这一概念的历史实质及其某种意义上的吊诡、双重含义首先表明，随着这一概念，从法国革命开始的市民社会与政治生活的分离尽管应当被扬弃，然而同时由英国开始的经济革命却要接受到这一概念之中。黑格尔市民社会的这种复杂的双重时代关联，其不仅指向过去和未来，而且指向现代革命的两股潮流，指向法国和英国，这一点构成了它的历史现实性，并且成为它在下一代发生政治影响的基础。

\section*{5}

即使我们并不接受里特尔新近出版的黑格尔解读\footnote{《黑格尔与法国大革命》，美茵河畔法兰克福，1965年新版。}的所有要点，那么市民社会在法哲学的概念结构中一再出现——作为抽象法的具体基础（第209节以下），作为主观道德的广泛的活动领域（在贫困领域，第242节），作为消失在其两极分化之中的“伦理事物的现象世界”以及作为伦理的“沦丧”（第182、184和229节），作为个体之“子”和家庭的实体性基础（第238节），最后作为国家的“前提”（第182节）——这一事实也就证明了他的这种论点的正确性：自1820年起，市民社会即“进入哲学及其政治理论的中心”。\footnote{《黑格尔与法国大革命》，第57页以下。}更加具体地说，市民社会尤其处于“伦理”的三个部分的中间，并且表现为将家庭和国家连接起来的部分，这在欧洲政治哲学传统中是前所未有的。由于这里不能对这个历史问题进行深究，因此就仅限于提到托马修斯在18世纪初如何界定市民社会。他在《政治明智短论》中写道：“然而人的社会本身要么是市民社会，要么是家庭社会。家庭社会构成市民社会的基础，因为这里市民社会的意思不外乎许多家庭社会及其人口的联合，而处于一个普遍的统治之下。”\footnote{托马修斯：《政治明智短论》，1705年，第204页以下。}这个定义完全处于旧传统之下，并且像《普鲁士国家的普通邦法》(Allgemeines Landrecht für die Preußischen Staaten, 1794)这样一部18世纪末制定的贴近当时时代社会现实的法典一样，当它在第一部分第2节宣称“市民社会由更多更小的、通过自然或法律（或是通过两者）结合起来的社会和等级组成”\footnote{《普鲁士国家的普通邦法》，第一部分，第1章，第2节。}时，也同样与传统相去不远。

与传统唯一的但却重大的区别也许是用诸等级来补充“社会”——沃尔夫的小型社会。在《普通邦法》中，市民社会仍然表现为政治社会、国家，其根据在于至少理论上保持不变的家庭、家庭社会的传统地位；相反，市民社会则以其政治公共领域的结构与之形成鲜明对照。因此，《普通邦法》下一节便立即规定：“夫妻之间、父母与子女之间的联合构成家庭社会。仆人也要算入家庭社会。”\footnote{《普鲁士国家的普通邦法》第一部分，第1章，第3节。}在这一点上，《普通邦法》的作者与同一时期康德的《道德形而上学》（1797年）第一部分即《法权论》完全一致。在康德的《法权论》中，市民社会同样保留下来，而家庭社会法也有其自己的部分。\footnote{《法权论》，第22节，见《康德全集》（科学院版）第6卷，第276页。在康德看来“家庭制度”的特征是“物的方式的人格权”，即“占有一个作为物的外部对象，并将这个对象作为一个人格来使用”（第276页）。通过这种法权关系，康德将家庭团体的古老“家政”结构、主奴关系（第283页）固定下来。对此，参见黑格尔在《法哲学原理》第40节对这种“物的方式的人格权”的批判。}然而，一旦“国民经济学”的发展突破了这种“家庭社会”的限制（如在研读过斯图亚特和斯密的黑格尔这里），社会便获得了指派给“oeconomia”（家政）的功能。相反，“societas civilis”作为“市民社会”便获得了新的地位，随着工业革命开始，市民社会的实体必然性支配了个人及其“小型社会”，而家庭社会则相应地解体。

我们可以在黑格尔的家庭概念上准确地看出这种交互影响。作为“伦理”的第一阶段，家庭不是一个拥有不同地位的成员的“社会”，它不是由丈夫和妻子、父母和子女的“小型社会”构成，而且绝对不会由主人和奴隶构成；相反，它是一个法律上的“人格”，“在所有物中”，而非在家庭中，有其外部实在性。\footnote{《法哲学原理》，第169节：“家庭作为人将在所有物中具有它的外部现实性。它只有在作为财富的所有物中才具有它的实体性的人格性的定在。”}不同于康德，在黑格尔这里，传统欧洲家庭的“经济”性消失了，并被18世纪晚期的“情感性”家庭概念所取代。\footnote{布伦纳在其论文《“完整家庭”和传统欧洲的“经济学”》(Das "Ganze Haus" und die alteuropäische "Ökonomik")中使用了这一表述，该文后被收入《社会史的新道路》（1956年），第44页。}《法哲学原理》第158节规定：“作为精神的直接实体性，家庭以精神的自我感觉到的统一性——爱——为规定，因此这里情感就是，精神的个体性在这种作为自在自为地存在着的本质性的统一性中自我意识到，自己在其中不是作为自为的个人，而是作为成员而存在。”

“必然性”作为市民社会的基本规定，消解了作为基本经济单位的家庭。我们看到，市民社会使其结构交织起来，结合为一个实体性的、经济的和法的整体。个人尽管还是市民社会的“成员”，然而首先不再是参与到其实体之中，而是作为“独立的”个人与市民社会的运动相对立。因为第238节说，唯有家庭才“首先”是“实体性整体”，它的职责在于“照料个人的特殊方面，既要考虑到他的手段和技能，使他能够从普遍财富中谋生，又要考虑到他在丧失工作能力时的生活和照顾”。然而实际上，如黑格尔所言，家庭“在市民社会中是从属的东西，它只构成基础：它的活动范围并不如此广泛。相反，市民社会才是惊人的力量，它把人扯到自己这里来，要求他为它工作，要求他的一切都要通过它，并且通过它而行动”。\footnote{《法哲学原理》，第238节，补充，第299页。（中译文参见《法哲学原理》，第241页。——译者）}这里，市民社会的必然性表现为它之所是：居于家庭的伦理实体和个人生活之上的“惊人的力量”，它使个人彼此“生疏”，只承认他们是“独立的个人”：“但是，市民社会把个人从这种联系中揪出，使家庭成员彼此生疏，并承认他们是独立的个人；此外，市民社会又用它自己的土地取代个人赖以为生的外部的无机自然和父辈的土地，使整个家庭的存在都依赖于它，听从偶然性的支配。由此，个人就成为市民社会的子女，市民社会要求于他，而他也对市民社会拥有权利。”\footnote{《法哲学原理》，第238节。（中译文参见《法哲学原理》，第241页。——译者）}这些地方最清楚地表达了黑格尔关于家庭团体中“小型社会”的解体及其与市民社会概念之间的联系的看法。

于是，黑格尔的市民社会概念涉及传统意义上的“家庭”社会的结构变化，同时另一方面也关系到“政治”国家的地位变化。这里，政治学不再是市民科学(scientia civilis)，即在其政治制度中对市民社会进行阐述，而是与之相对的“国家科学”。在欧洲旧传统中，国家和社会在市民社会的关系概念中交织在一起，在黑格尔这里，它们才开始被“设定”到“关系”之中，并且彼此处于一种特殊的“差别”之中：“市民社会是处于家庭和国家之间的差别[阶段]，尽管其形成要晚于国家。因为它作为差别，以国家为前提，它要存在，前面就必须要有一个独立的国家。”\footnote{同上书，第182节，补充，第246页。（中译文参见《法哲学原理》，第196页。——译者）}黑格尔将现代的国家制度的现实性纳入神一般静止的、实现了其目的的、自足的古代城邦理念的传统形式之中，在其内部解决了国家与社会之间的二元论。因为黑格尔明言，这个国家乃是“现代国家”，其本质在于，“普遍物是同特殊性的完全自由和私人福利相结合的，所以家庭和市民社会的利益必须集中于国家；但是，特殊性也必须保持其权利，如果没有特殊性自己的知识和意志，那么目的的普遍性也就不能前进”。\footnote{同上书，第260节，补充，第322页。（中译文参见《法哲学原理》，第261页。——译者）}

在黑格尔这里，国家不仅“设定”了市民社会，而且相反，还为其自身的制度现实性预设了市民社会的存在。柏林历史哲学讲座对北美政治状况所作的令人印象深刻的经验分析表明了这一点。这里，黑格尔得出了这样的结论：现在北美仍然不是一个完全发展了的国家，因为它还没有市民社会。在北美，这种发展受到西部土地占领的阻碍，人口的“不断迁移”阻碍了城镇化进程：“只有像在欧洲一样，农业的简单增长受阻，居民才会不再走向田地，而是涌向城市工商业，形成紧密的市民社会体系，满足有机国家的需要。”\footnote{《历史哲学》，导论，引自《历史中的理性》，1955年，第209页。}这里，如布伦奇利在19世纪晚期所言，市民社会已经带有了“第三等级概念”的特征。

黑格尔与传统的断裂并不只是表现在这一点上面，即政治学的对象分化为国家和市民社会的差别。它更是通过这种现代“国家”与历史的关系而表现出来。因为它的前提并不仅仅基于市民社会的动力学，而是同时向上延伸到历史之流中。国家理念在其臻于完成的顶点表现为“类以及与个别国家相对的绝对权力、在世界历史过程中赋予自己现实性的精神”。\footnote{《法哲学原理》，第259节。（中译文参见《法哲学原理》，第259页。——译者）}国家不再是传统的政治状态概念，不再是贯穿传统政治思想之始终的国家或市民社会的静止不变的自然模型；相反，历史作为新的上司高踞其上。在历史中，“伦理的整体自身、国家的独立性”——这里我们必须替黑格尔补充完整：再次——“委诸偶然性之下”。\footnote{同上书，第340节。（中译文参见《法哲学原理》，第351页。——译者）}对于黑格尔《法哲学原理》的这种结尾，甘斯言道，“真正说来”，其中包含了“此书最重要的价值”。一般而言，这种说法道出了他那一代人的心声和19世纪观念。在甘斯看来，这个结尾乃是一幅“壮丽的景象”，从国家的高度出发，我们看到众多国家“卷入历史的世界海洋”；同时，他在当时率先指出，黑格尔在《法哲学原理》末尾给出的“历史发展概要只是对这一领域中那些更加重要的关切的一个预告”。\footnote{《法哲学原理》，前言，第VIII-IX页。}

\section*{6}

如果我们全面考察黑格尔法哲学的概念结构及其建构原则，那就会说，作为革命的结果和遗产，市民社会和历史这两个要素耸立其中。两者同时改变了自古代以来一直有效的政治哲学传统所阐述的家庭和国家的内部制度。正如家庭居于作为家庭成员的个人与市民社会“之间”一样，国家运行于市民社会与历史之间。在黑格尔这里，正是在古代实践哲学体系的建筑术上的起源，产生了法的富于艺术性的辩证结构。最后，法哲学作为“中介体系”由此出发获得了富有历史内容的独特财富，这种财富使它在原则上与卢梭的《社会契约论》、康德的《法权论》和费希特的《自然法基础》等政治著作区分开来。

黑格尔在古代哲学史讲座中言及柏拉图时，反对对《理想国》的流行成见，说柏拉图“不是一个玩弄抽象理论和抽象原则的人，他的真实精神曾经认识了并表述了真实的事物。这不能是别的，而只能是他生活于其中的世界的真实事物，也只能是那唯一很好地活在他本人和希腊里面的[时代]精神的真实事物”。在更深的意义上，这种说法也适用于他本人。事实上，他用《法哲学原理》以现代世界的“实体性方式”描述了这个世界的真正“伦理”。他用来结束柏拉图讲座的这一部分的那些话语，也能很好地放在作为最后的形而上学的哲学“政治学”的法哲学的开头：“没有人能够跳出他的时代，他的时代的精神也就是他的精神；但问题在于认识到时代精神的具体内容。”\footnote{《哲学史讲演录》，《全集》第14卷，第275页。（中译文引自《哲学史讲演录》第二卷，第248页。——译者）}